\chapter{Heisenberg Uncertainty Principle}
We now present the mathematical foundation of uncertainty relations in quantum mechanics. \\
Unlike classical mechanics, where a particle's position and momentum can be simultaneously determined with arbitrary precision, quantum mechanics imposes fundamental limitations on measurement precision. We can only specify:

\begin{align}
& \left\langle x_{j}\right\rangle \pm \Delta x_{j} \\
& \left\langle\hat{p}_{j}\right\rangle \pm \Delta \hat{p}_{j}
\end{align}

This inherent uncertainty is not due to technological limitations or experimental errors but represents a fundamental property of nature. Our objective is to derive the general Heisenberg uncertainty principle:

\begin{equation}
\left(\Delta \hat{A}_{1}\right)^{2}\left(\Delta \hat{A}_{2}\right)^{2} \geq \frac{1}{4}\left|\left\langle\left[\hat{A}_{1}, \hat{A}_{2}\right]\right\rangle\right|^{2}
\end{equation}


\section{Derivation of the General Uncertainty Principle}

Let's consider two physical observables represented by Hermitian operators $\hat{A}_1$ and $\hat{A}_2$:

\begin{align}
& \hat{A}_{1}^{\dagger}=\hat{A}_{1} \\
& \hat{A}_{2}^{\dagger}=\hat{A}_{2}
\end{align}

We define the uncertainty (or standard deviation) of a physical quantity as:

\begin{equation}
(\Delta \hat{A})^{2}=\left(\Psi,(\hat{A}-\langle\hat{A}\rangle)^{2} \Psi\right)
\end{equation}

For convenience, we introduce deviation operators:

\begin{align}
& \hat{D}_{1}=\hat{A}_{1}-\left\langle\hat{A}_{1}\right\rangle \\
& \hat{D}_{2}=\hat{A}_{2}-\left\langle\hat{A}_{2}\right\rangle
\end{align}

The uncertainties can then be expressed as:

\begin{align}
& \left(\Delta \hat{A}_{1}\right)^{2}=\left(\Psi, \hat{D}_{1}^{2} \Psi\right) \\
& \left(\Delta \hat{A}_{2}\right)^{2}=\left(\Psi, \hat{D}_{2}^{2} \Psi\right)
\end{align}

Expanding $\hat{D}_j^2$:

\begin{equation}
\hat{D}_{j}^{2}=\left(\hat{A}_{j}-\left\langle\hat{A}_{j}\right\rangle\right)^{2}=\hat{A}_{j}^{2}-2\left\langle\hat{A}_{j}\right\rangle \hat{A}_{j}+\left\langle\hat{A}_{j}\right\rangle^{2}
\end{equation}

Taking the expectation value:

\begin{equation}
\left\langle\hat{D}_{j}^{2}\right\rangle=\left\langle\hat{A}_{j}^{2}\right\rangle-2\left\langle\hat{A}_{j}\right\rangle\left\langle\hat{A}_{j}\right\rangle+\left\langle\hat{A}_{j}\right\rangle^{2}=\left\langle\hat{A}_{j}^{2}\right\rangle-\left\langle\hat{A}_{j}\right\rangle^{2}=\sigma_{A_{j}}^{2}
\end{equation}

The commutator of the deviation operators equals the commutator of the original operators:

\begin{align}
{\left[\hat{D}_{1}, \hat{D}_{2}\right] } & =\left[\hat{A}_{1}-\left\langle\hat{A}_{1}\right\rangle, \hat{A}_{2}-\left\langle\hat{A}_{2}\right\rangle\right] \\
& =\left[\hat{A}_{1}, \hat{A}_{2}-\left\langle\hat{A}_{2}\right\rangle\right]-\left[\left\langle\hat{A}_{1}\right\rangle, \hat{A}_{2}-\left\langle\hat{A}_{2}\right\rangle\right] \\
& =\left[\hat{A}_{1}, \hat{A}_{2}\right]-\left[\hat{A}_{1},\left\langle\hat{A}_{2}\right\rangle\right]-\left[\left\langle\hat{A}_{1}\right\rangle, \hat{A}_{2}\right]+\left[\left\langle\hat{A}_{1}\right\rangle,\left\langle\hat{A}_{2}\right\rangle\right]  \\
& =\left[\hat{A}_{1}, \hat{A}_{2}\right]
\end{align}

Since $\hat{D}_1$ and $\hat{D}_2$ are Hermitian, we can establish:

\begin{align}
\left\langle\hat{D}_{1} \hat{D}_{2}\right\rangle^{*} & =\left(\Psi, \hat{D}_{1} \hat{D}_{2} \Psi\right)^{*}= \\
& =\left(\hat{D}_{1} \hat{D}_{2} \Psi, \Psi\right)=  \\
& =\left(\hat{D}_{2} \Psi, \hat{D}_{1} \Psi\right)= \\
& =\left(\Psi, \hat{D}_{2} \hat{D}_{1} \Psi\right)=\left\langle\hat{D}_{2} \hat{D}_{1}\right\rangle
\end{align}

This gives us:

\begin{equation}
\left\langle\hat{D}_{1} \hat{D}_{2}\right\rangle^{*}=\left\langle\hat{D}_{2} \hat{D}_{1}\right\rangle
\end{equation}

Using these relations, we can evaluate:

\begin{align}
\left|\left\langle\left[\hat{A}_{1}, \hat{A}_{2}\right]\right\rangle\right|^{2} & \stackrel{(1)}{=}\left|\left\langle\left[\hat{D}_{1}, \hat{D}_{2}\right]\right\rangle\right|^{2}=\left|\left\langle\hat{D}_{1} \hat{D}_{2}-\hat{D}_{2} \hat{D}_{1}\right\rangle\right|^{2}= \\
& =\left|\left\langle\hat{D}_{1} \hat{D}_{2}\right\rangle-\left\langle\hat{D}_{1} \hat{D}_{2}\right\rangle^{*}\right|^{2}=\left|2 i \operatorname{Im}\left(\left\langle\hat{D}_{1} \hat{D}_{2}\right\rangle\right)\right|^{2}= \\
& =4\left|\operatorname{Im}\left(\left\langle\hat{D}_{1} \hat{D}_{2}\right\rangle\right)\right|^{2} \leq 4\left|\operatorname{Re}\left(\left\langle\hat{D}_{1} \hat{D}_{2}\right\rangle\right)\right|^{2}+4\left|\operatorname{Im}\left(\left\langle\hat{D}_{1} \hat{D}_{2}\right\rangle\right)\right|^{2}= \\
& =4\left|\left\langle\hat{D}_{1} \hat{D}_{2}\right\rangle\right|^{2}
\end{align}

This can be rewritten as:

\begin{equation}
4\left|\left\langle\hat{D}_{1} \hat{D}_{2}\right\rangle\right|^{2}=4\left(\Psi,\left(\hat{D}_{1} \hat{D}_{2}\right) \Psi\right)^{2}=4\left(\hat{D}_{1} \Psi, \hat{D}_{2} \Psi\right)^{2}
\end{equation}

Defining $\hat{D}_{1} \Psi=\phi_{1}$ and $\hat{D}_{2} \Psi=\phi_{2}$, we can apply the Cauchy-Schwarz inequality:

\begin{equation}
\left(\phi_{1}, \phi_{2}\right)^{2} \leq\left|\phi_{1}\right|^2\left|\phi_{2}\right|^2
\end{equation}


\section{Deriving the Uncertainty Relation}
This analysis leads to:

\begin{align}
& 4\left(\Psi,\left(\hat{D}_{1} \hat{D}_{2}\right) \Psi\right)^{2} \stackrel{(2)}{\leq} 4\left|\hat{D}_{1} \Psi\right|\left|\hat{D}_{2} \Psi\right|= \\
& =4\left(\hat{D}_{1} \Psi, \hat{D}_{1} \Psi\right)\left(\hat{D}_{2} \Psi \hat{D}_{2} \Psi\right)=4\left(\Psi, \hat{D}_{1}^{2} \Psi\right)\left(\Psi \hat{D}_{2}^{2} \Psi\right)  \\
& =4\left\langle\hat{D}_{1}^{2}\right\rangle\left\langle\hat{D}_{2}^{2}\right\rangle=4\left(\Delta \hat{A}_{1}\right)^{2}\left(\Delta \hat{A}_{2}\right)^{2} \stackrel{(3)}{=} 4 \sigma_{A_{1}}^{2} \sigma_{A_{2}}^{2}
\end{align}

Combining the previous results yields:

\begin{equation}
\left|\left\langle\left[\hat{A}_{1}, \hat{A}_{2}\right]\right\rangle\right|^{2}=\left|\left\langle\left[\hat{D}_{1}, \hat{D}_{2}\right]\right\rangle\right|^{2} \leq 4\left|\hat{D}_{1} \Psi\right|\left|\hat{D}_{2} \Psi\right|=4 \sigma_{A_{1}}^{2} \sigma_{A_{2}}^{2}
\end{equation}

Thus we arrive at the fundamental uncertainty relation:

\begin{equation}
\sigma_{A_{1}}^{2} \sigma_{A_{2}}^{2} \geq \frac{1}{4}\left|\left\langle\left[\hat{A}_{1}, \hat{A}_{2}\right]\right\rangle\right|^{2}
\end{equation}

\section{Physical Implications}
The Heisenberg uncertainty principle reveals a fundamental limitation: increasing measurement precision for one observable necessarily reduces precision for another non-commuting observable. The only exception occurs when operators commute, yielding $\sigma_{A_{1}}^{2} \sigma_{A_{2}}^{2} \geq 0$, allowing simultaneous precise measurements.

The classical notion of simultaneously determining a particle's position and momentum with arbitrary precision fails in quantum mechanics. From the commutation relation:

\begin{equation}
\left[x_{i}, p_{j}\right]=i \hbar \delta_{i j}
\end{equation}

We derive:

\begin{equation}
\sigma_{x_{i}}^{2} \sigma_{p_{j}}^{2} \geq \frac{1}{4}\left|\left\langle\left[x_{i}, p_{j}\right]\right\rangle\right|^{2}=\frac{\hbar^{2}}{4} \delta_{i j}
\end{equation}

For the same coordinate direction ($i=j$), we obtain the iconic inequality:

\begin{equation}
\Delta x \Delta p \geq \frac{\hbar}{2}
\end{equation}

\section{Properties of Commuting Operators}
For a wavefunction characterized by two quantum numbers $\alpha_1$ and $\alpha_2$, the following eigenvalue equations hold simultaneously if and only if the operators commute:

\begin{align}
& \hat{A}_{1} \Psi_{\alpha_{1}, \alpha_{2}}=\alpha_{1} \Psi_{\alpha_{1}, \alpha_{2}} \\
& \hat{A}_{2} \Psi_{\alpha_{1}, \alpha_{2}}=\alpha_{2} \Psi_{\alpha_{1}, \alpha_{2}}
\end{align}

This follows from:

\begin{equation}
\left[\hat{A}_{1}, \hat{A}_{2}\right]=0 \Longleftrightarrow \hat{A}_{1} \hat{A}_{2} \Psi_{\alpha_{1}, \alpha_{2}}-\hat{A}_{2} \hat{A}_{1} \Psi_{\alpha_{1}, \alpha_{2}}
\end{equation}

For this relation to hold, $\hat{A}_1$ and $\hat{A}_2$ must have eigenvalues $\alpha_1$ and $\alpha_2$.

Theorem 7.1. If two operators $\hat{A}$ and $\hat{B}$ satisfy $[\hat{A}, \hat{B}]=0$, they possess a common basis of eigenfunctions.

\section{Non-degenerate case}
Proof. For an eigenstate $\Psi$ of $\hat{A}$ with $\hat{A} \Psi=\alpha \Psi$, consider $\Psi_{B}=\hat{B} \Psi$. Then:

\begin{equation}
\hat{A} \Psi_{B}=\hat{A} \hat{B} \Psi=\hat{B} \hat{A} \Psi=\alpha \hat{B} \Psi=\alpha \Psi_{B}
\end{equation}

Thus $\Psi_{B}$ is also an eigenstate of $\hat{A}$. Since $\Psi_B$ is an eigenstate, we must have $\hat{A} \Psi_{B}=\gamma \Psi_{B}$. This requires:

\begin{equation}
\hat{B} \Psi=\beta \Psi \Longrightarrow \gamma=\alpha \beta
\end{equation}

Therefore, $\Psi$ must be a simultaneous eigenstate of both $\hat{A}$ and $\hat{B}$.

\section{Degenerate case}
Proof. For the degenerate case, we consider the basis:

\begin{equation}
\mathcal{B}=\left\{\Psi_{i}: i \in[1, M], \hat{A} \Psi_{i}=\alpha_{i} \Psi_{i}\right\}
\end{equation}


\section{Analyzing Commuting Operators}
Continuing with the degenerate case:

\begin{equation}
\hat{A} \hat{B} \Psi_{i}=\hat{B} \hat{A} \Psi_{i}=\alpha_{i} \hat{B} \Psi
\end{equation}

While $\hat{B} \Psi_{i}$ represents a wavefunction, generally $\hat{B} \Psi_{i} \neq \beta \Psi_{i}$. Instead, $\hat{B} \Psi_{i}$ can be expressed as a linear combination:

\begin{equation}
\hat{B} \Psi_{i}=\sum_{n=1}^{M} C_{i n} \Psi_{n}
\end{equation}

Where $C_{i n}=\left(\Psi_{i}, \hat{B} \Psi_{n}\right)$. Examining the Hermiticity of $C_{in}$:

\begin{equation}
C_{i n}^{\dagger}=C_{n i}^{*}=\left(\Psi_{n}, \hat{B} \Psi_{i}\right)^{*}=\left(\hat{B} \Psi_{i}, \Psi_{n}\right)=\left(\Psi_{i}, \hat{B}^{\dagger} \Psi_{n}\right)=\left(\Psi_{i}, \hat{B} \Psi_{n}\right)=C_{i n}
\end{equation}

Thus $C_{i n}$ forms a Hermitian matrix. For any Hermitian matrix $C$, there exists a unitary matrix $U$ such that:

\begin{align}
& U^{\dagger} C U=D \quad \text { where } D \text { is a diagonal matrix } \\
& U^{\dagger} U=\mathbb{I}
\end{align}

In component form:

\begin{equation}
U_{i j}^{\dagger} C_{j n} U_{n m}=U_{j i}^{*} C_{j n} U_{n m}=d_{i} \delta_{i j}
\end{equation}

With the unitarity condition:

\begin{equation}
U_{i k} U_{k j}^{\dagger}=\delta_{i j}
\end{equation}

Defining new states $\phi_{i}=U_{j i}^{*} \Psi_{j}$, we prove these are eigenstates of $\hat{B}$:

\begin{equation}
\hat{B} \phi_{i}=\hat{B} U_{j i}^{*} \Psi_{j}=U_{j i}^{*} \hat{B} \Psi_{j}=U_{j i}^{*} C_{j n} \Psi_{n}=U_{j i}^{*} C_{j n} \delta_{n m} \Psi_{m}
\end{equation}

Substituting the identity matrix with our $U$ matrix definition:

\begin{equation}
U_{j i}^{*} C_{j n} \delta_{n m} \Psi_{m}=\underbrace{U_{j i}^{*} C_{j n} U_{n k}}_{d_{i} \delta_{i k}} U_{k m}^{\dagger} \Psi_{m}
\end{equation}

This yields:

\begin{equation}
d_{i} \delta_{i k} U_{k m}^{\dagger} \Psi_{m}=d_{i} U_{m i}^{*} \Psi_{m}=d_{i} \phi_{i}
\end{equation}

Confirming $\phi_{i}$ is an eigenstate of $\hat{B}$ with eigenvalue $d_i$. Additionally:

\begin{equation}
\hat{A} \phi_{i}=\hat{A} U_{j i}^{*} \Psi_{j}=U_{j i}^{*} \hat{A} \Psi_{j}=U_{j i}^{*} \alpha \Psi_{j}=\alpha \phi_{i}
\end{equation}

Thus $\phi_{i}$ is a simultaneous eigenstate of both $\hat{A}$ and $\hat{B}$.

\section{Key Implications}
The commutation relation $[\hat{A}, \hat{B}]=0$ yields several significant consequences:

\begin{enumerate}
  \item If $[\hat{A}, \hat{H}]=0$, applying the Ehrenfest theorem we have:
\end{enumerate}


\section{Consequences of Operator Commutation}

\begin{equation}
i \hbar \frac{\, \mathrm{d}\langle\hat{A}\rangle}{\mathrm{d} t}=\langle[\hat{A}, \hat{H}]\rangle=0 \Longrightarrow \frac{\, \mathrm{d}\langle\hat{A}\rangle}{\mathrm{d} t}=0
\end{equation}

This demonstrates that when an operator commutes with the Hamiltonian, its expectation value remains constant over time.

Additional consequences include:
\begin{enumerate}
  \item If $[\hat{A}, \hat{H}]=0$, a common eigenbasis exists
  \item For commuting operators, the uncertainty relation reduces to:

\begin{equation}
(\Delta \hat{A})^{2}(\Delta \hat{H})^{2} \geq 0
\end{equation}

  Allowing for states where both $\Delta \hat{A}=0$ and $\Delta \hat{H}=0$
  \item Complete sets of observables can be defined
  \item Coherent states become possible
\end{enumerate}

We now examine the last two points in detail.

\section{Complete Observable Sets}
A collection of Hermitian operators $\hat{A}, \hat{B}, \ldots, \hat{P}$ forms a complete set when it contains the maximum possible number of mutually commuting operators.

Example 1: For a free particle Hamiltonian:

\begin{equation}
\hat{H}=\frac{\hat{p}_{1}^{2}+\hat{p}_{2}^{2}+\hat{p}_{3}^{2}}{2 m}
\end{equation}

Several operator sets can be considered:
\begin{enumerate}
  \item $x_{1}, x_{2}, x_{3}$ - mutually commuting but not with $\hat{H}$
  \item $\hat{p}_{1}, \hat{p}_{2}, \hat{p}_{3}$ - mutually commuting and commuting with $\hat{H}$
  \item $x_{1}, x_{2}, p_{3}$ - mutually commuting but not with $\hat{H}$
  \item $\hat{L}_{3}, \hat{L}^{2}, \hat{H}$ - mutually commuting
  \item $\hat{p}_{i}, \hat{L}_{i}, \hat{H}$ - mutually commuting
\end{enumerate}

While options 4 and 5 are valid, $\hat{p}_{1}, \hat{p}_{2}, \hat{p}_{3}$ is most practical. For plane waves:

\begin{equation}
\hat{p}_{j} \phi_{\vec{k}}(\vec{x})=\hbar k_{j} \phi_{\vec{k}}(\vec{x})
\end{equation}

The Hamiltonian acts on these eigenfunctions as:

\begin{equation}
\hat{H} \phi_{\vec{k}}(\vec{x})=\frac{1}{2 m} \sum_{j} \hat{p}_{j}^{2} \phi_{\vec{k}}(\vec{x})=\frac{\hbar k^{2}}{2 m} \phi_{\vec{k}}(\vec{x})
\end{equation}

Example 2: For a Hamiltonian with potential:

\begin{equation}
\hat{H}=\frac{\hat{p}^{2}}{2 m}+U(\vec{x})
\end{equation}

Many commuting sets exist:
\begin{enumerate}
  \item $x_{1}, x_{2}, x_{3}$
  \item $\hat{p}_{1}, \hat{p}_{2}, \hat{p}_{3}$\\
$\cdots$\\
n. $\hat{L}_{3}, \hat{L}^{2}, \hat{H}$
\end{enumerate}

Without knowledge of the potential form, any set might be appropriate.

Example 3: For a spherically symmetric potential:

\begin{equation}
\hat{H}=\frac{\hat{p}^{2}}{2 m}+U(|\vec{x}|)
\end{equation}

The set $\hat{L}_{3}, \hat{L}^{2}, \hat{H}$ becomes particularly advantageous due to rotational symmetry. Note that since $\left[\hat{L}_{i}, \hat{L}_{j}\right]=i \hbar \epsilon_{i j k} \hat{L}_{k}$, we cannot simultaneously include both $\hat{L}_{2}$ and $\hat{L}_{3}$, but must use $\hat{L}^{2}$.

\section{Quantum Coherent States}
Coherent states minimize the uncertainty relation:

\begin{equation}
\Delta \hat{A} \Delta \hat{B}=\frac{1}{2}|[\hat{A}, \hat{B}]|
\end{equation}

These states exhibit quasi-classical behavior with minimal observable uncertainties. For position and momentum:

\begin{equation}
\Delta x \Delta \hat{p}=\frac{\hbar}{2}
\end{equation}


\section{Analysis of Quantum Coherent States}
Coherent states satisfy the uncertainty relation with the minimal possible product:

\begin{equation}
\Delta \hat{A}=\sqrt{\left\langle\hat{A}^{2}\right\rangle-\langle\hat{A}\rangle^{2}}
\end{equation}

For position and momentum, coherent states are defined as eigenstates of the annihilation operator:

\begin{equation}
\hat{a} \Psi_{z}=z \Psi_{z}
\end{equation}

Using the definition of the annihilation operator:

\begin{equation}
\hat{a}=\frac{1}{\sqrt{2} \lambda}\left(x+i \frac{\lambda^{2} \hat{p}}{\hbar}\right)
\end{equation}

The eigenvalue equation becomes:

\begin{equation}
\frac{1}{\sqrt{2} \lambda}\left(x+\lambda^{2} \frac{\partial}{\partial x}\right) \Psi_{z}=z \Psi_{z}
\end{equation}

This first-order differential equation can be solved by proposing:

\begin{equation}
\Psi_{z}=c \mathrm{e}^{-\sigma x^{2}+i \eta x}
\end{equation}

Substituting into the eigenvalue equation:

\begin{align}
& \frac{1}{\sqrt{2} \lambda}\left[x+\lambda^{2} \frac{\partial}{\partial x}\right]\left(c e^{-\sigma x^{2}+i \eta x}\right)= \\
& =\frac{1}{\sqrt{2} \lambda}\left[x c e^{-\sigma x^{2}+i \eta x}+\lambda^{2} \frac{\partial}{\partial x}\left(c e^{-\sigma x^{2}+i \eta x}\right)\right]= \\
& =\frac{1}{\sqrt{2} \lambda}\left[x c e^{-\sigma x^{2}+i \eta x}+\lambda^{2} c(-2 \sigma x+i \eta) e^{-\sigma x^{2}+i \eta x}\right]  \\
& =\frac{c e^{\sigma x^{2}+i \eta x}}{\sqrt{2} \lambda}\left[x+\lambda^{2}(-2 \sigma x+i \eta)\right]
\end{align}

Which gives:

\begin{equation}
z c e^{-\sigma x^{2}+i \eta x}=\frac{c e^{-\sigma x^{2}+i \eta x}}{\sqrt{2} \lambda}\left[x-\lambda^{2}(2 \sigma x+i \eta)\right]
\end{equation}

For this equation to be satisfied:

\begin{align}
& 1-\lambda^{2} 2 \sigma=0 \Longrightarrow \sigma=\frac{1}{2 \lambda^{2}} \\
& z=\frac{i \lambda \eta}{\sqrt{2}} \Longrightarrow i \eta=\frac{\sqrt{2}}{\lambda} z
\end{align}

Therefore, the coherent state wavefunction is:

\begin{equation}
\Psi_{z}=c \exp \left(-\frac{x^{2}}{2 \lambda^{2}}+\frac{\sqrt{2}}{\lambda} z x\right) \label{eq:7.53}
\end{equation}

To normalize this wavefunction, we define $\nu=\operatorname{Re}(z)$ so that $z+z^{*}=2 \nu$ and calculate:

\begin{equation}
\Psi_{z} \Psi_{z}^{*}=c^{2} \exp \left(-\frac{x^{2}}{\lambda^{2}}+\frac{\sqrt{2}}{\lambda}\left(z+z^{*}\right) x\right)=c^{2} \exp \left(-\frac{x^{2}}{\lambda^{2}}+\frac{\sqrt{2}}{\lambda} \nu x\right)
\end{equation}

Completing the square in the exponent:

\begin{equation}
x^{2}-2 \sqrt{2} \lambda \nu x=(x-\sqrt{2} \lambda \nu)^{2}-2 \lambda^{2} \nu^{2}
\end{equation}

The probability density becomes:

\begin{equation}
\Psi_{z} \Psi_{z}^{*}=c^{2} \exp \left(-\frac{1}{\lambda^{2}}(x-\sqrt{2} \lambda \nu)^{2}+2 \nu^{2}\right)
\end{equation}

This represents a Gaussian function centered at $x_{0}=\frac{\lambda \nu}{\sqrt{2}}$.

The normalization integral is:

\begin{align}
& c^{2} \int_{-\infty}^{+\infty} \mathrm{d} x \exp \left(-\frac{1}{\lambda^{2}}(x-\sqrt{2} \lambda \nu)^{2}+2 \nu^{2}\right)= \\
& c^{2} \exp \left(2 \nu^{2}\right) \int_{-\infty}^{+\infty} \mathrm{d} x \exp \left(-\frac{1}{\lambda^{2}}(x-\sqrt{2} \lambda \nu)^{2}\right)
\end{align}

I understand now - you want me to revise the LaTeX code without adding comments, while maintaining mathematical expressions exactly as they are. Here's the revised code:


This represents a gaussian integral with parameter $a=\frac{1}{\lambda^{2}}$ yielding $\sqrt{\frac{\pi}{a}}=\lambda \sqrt{\pi}$. The normalization condition requires:

\begin{equation}
c^{2} \mathrm{e}^{2 \nu^{2}} \lambda \sqrt{\pi}=1 \Longrightarrow c=\frac{\mathrm{e}^{-\nu^{2}}}{\sqrt{\lambda \sqrt{\pi}}}
\end{equation}

With the wavefunction normalized, we can determine the explicit form of $\Psi_{z}$. We know:

\begin{equation}
\left(\Psi_{z}, \hat{a} \Psi_{z}\right)=z\left(\Psi_{z}, \Psi_{z}\right)=z
\end{equation}

Additionally:

\begin{align}
\left(\Psi_{z}, \hat{a} \Psi_{z}\right) & =\frac{1}{\sqrt{2} \lambda}\left(\Psi_{z}, x \Psi_{z}\right)+\frac{i \lambda}{\sqrt{2} \hbar}\left(\Psi_{z}, \hat{p} \Psi_{z}\right)  \\
& =\frac{1}{\sqrt{2} \lambda}\langle x\rangle+\frac{i \lambda}{\sqrt{2} \hbar}\langle\hat{p}\rangle
\end{align}

Substituting into equation \eqref{eq:7.53}:

\begin{align}
\Psi_{z} & =\frac{\mathrm{e}^{-\nu^{2}}}{\sqrt{\lambda \sqrt{\pi}}} \exp \left(-\frac{x^{2}}{2 \lambda^{2}}+\frac{\sqrt{2}}{\lambda} \frac{1}{\sqrt{2} \lambda}\langle x\rangle x+\frac{\sqrt{2}}{\lambda} \frac{i \lambda}{\sqrt{2} \hbar}\langle\hat{p}\rangle x\right)= \\
& =\frac{\mathrm{e}^{-\nu^{2}}}{\sqrt{\lambda \sqrt{\pi}}} \exp \left(-\frac{x^{2}}{2 \lambda^{2}}+\frac{1}{\lambda^{2}}\langle x\rangle x+\frac{i}{\hbar}\langle\hat{p}\rangle x\right)= \\
& =\frac{\mathrm{e}^{-\nu^{2}}}{\sqrt{\lambda \sqrt{\pi}}} \exp \left[-\frac{1}{2 \lambda^{2}}\left(x^{2}-2\langle x\rangle x\right)+\frac{i}{\hbar}\langle\hat{p}\rangle x\right]=  \label{eq:7.61}\\
& =\frac{\mathrm{e}^{-\nu^{2}}}{\sqrt{\lambda \sqrt{\pi}}} \exp \left[-\frac{1}{2 \lambda^{2}}(x-\langle x\rangle)^{2}+\frac{\langle x\rangle^{2}}{2 \lambda^{2}}+\frac{i}{\hbar}\langle\hat{p}\rangle x\right]
\end{align}

For $\langle x\rangle$, the ladder operator properties give us:

\begin{equation}
x=\frac{\lambda}{\sqrt{2}}\left(\hat{a}+\hat{a}^{\dagger}\right)
\end{equation}

Since $z^{*}$ corresponds to $\hat{a}^{\dagger}$:

\begin{equation}
\left(\Psi_{z}, \hat{a}^{\dagger} \Psi_{z}\right)^{*}=\left(\hat{a}^{\dagger} \Psi_{z}, \Psi_{z}\right)=\left(\Psi_{z}, \hat{a} \Psi_{z}\right)=z
\end{equation}

Therefore:

\begin{align}
\langle x\rangle & =\left(\Psi_{z}, \frac{\lambda}{\sqrt{2}}\left(\hat{a}+\hat{a}^{\dagger}\right) \Psi_{z}\right)=\frac{\lambda}{\sqrt{2}}\left[\left(\Psi_{z}, \hat{a} \Psi_{z}\right)+\left(\Psi_{z}, \hat{a}^{\dagger} \Psi_{z}\right)\right]=  \\
& =\frac{\lambda}{\sqrt{2}}\left(z+z^{*}\right)=\frac{2 \nu \lambda}{\sqrt{2}}
\end{align}

Squaring this result:

\begin{equation}
\langle x\rangle^{2}=2 \nu^{2} \lambda^{2}
\end{equation}

Substituting into equation \eqref{eq:7.61}:

\begin{gather}
\Psi_{z}=\frac{1}{\sqrt{\lambda \sqrt{\pi}}} \exp \left[-\frac{1}{2 \lambda^{2}}(x-\langle x\rangle)^{2}-\mathcal{\nu}^{2}+\mathcal{P}^{2}+\frac{i}{\hbar}\langle\hat{p}\rangle x\right]  \\
\Psi_{z}=\frac{1}{\sqrt{\lambda \sqrt{\pi}}} \exp \left(-\frac{1}{2 \lambda^{2}}(x-\langle x\rangle)^{2}\right) \exp \left(i \frac{\langle\hat{p}\rangle}{\hbar} x\right)
\end{gather}

The second exponential represents a plane wave, while the first describes a gaussian distribution centered at $\langle x\rangle$ with variance $\sigma^{2}=\lambda^{2}$. This confines the highest probability region to a narrow band proportional to $\lambda$, explaining why coherent states approach deterministic behavior.

For the uncertainty relation $\Delta x \Delta p=\hbar / 2$, recall from the harmonic oscillator:

\begin{align}
\Delta x & =\sqrt{\left\langle x^{2}\right\rangle-\langle x\rangle^{2}}=\sqrt{\left\langle x^{2}\right\rangle} \\
\Delta p & =\sqrt{\left\langle\hat{p}^{2}\right\rangle-\langle\hat{p}\rangle^{2}}=\sqrt{\left\langle p^{2}\right\rangle}
\end{align}

\begin{equation}
\Delta x \Delta p=\frac{\hbar}{2}(2 n+1) \geq \frac{\hbar}{2}
\end{equation}

This relation represents the fundamental uncertainty principle for quantum states, establishing the minimum bound for simultaneous measurement precision.

\begin{align}
\langle x\rangle & =\sqrt{\frac{\hbar}{2 m \omega}}\left(\Psi_{z},\left(\hat{a}+\hat{a}^{\dagger}\right) \Psi_{z}\right)= \\
& =\sqrt{\frac{\hbar}{2 m \omega}}\left[\left(\Psi_{z}, \hat{a} \Psi_{z}\right)+\left(\Psi_{z}, \hat{a}^{\dagger} \Psi_{z}\right)\right]=  \\
& =\sqrt{\frac{\hbar}{2 m \omega}}\left(z+z^{*}\right)
\end{align}

Here we derive the expectation value of position for a coherent state by expressing the position operator in terms of creation and annihilation operators.

\begin{align}
\langle\hat{p}\rangle & =-i \sqrt{\frac{\hbar m \omega}{2}}\left(\Psi_{z},\left(\hat{a}-\hat{a}^{\dagger}\right) \Psi_{z}\right)= \\
& =-i \sqrt{\frac{\hbar m \omega}{2}}\left[\left(\Psi_{z}, \hat{a} \Psi_{z}\right)-\left(\Psi_{z}, \hat{a}^{\dagger} \Psi_{z}\right)\right]=  \\
& =-i \sqrt{\frac{\hbar m \omega}{2}}\left(z-z^{*}\right) \\
\left\langle x^{2}\right\rangle & =\frac{\hbar}{2 m \omega}\left(\Psi_{z},\left(\hat{a}+\hat{a}^{\dagger}\right)^{2} \Psi_{z}\right) \\
& =\frac{\hbar}{2 m \omega}\left(\Psi_{z},\left(\hat{a}^{2}+\hat{a}^{\dagger 2}+\hat{a} \hat{a}^{\dagger}+\hat{a}^{\dagger} \hat{a}\right) \Psi_{z}\right) \\
& =\frac{\hbar}{2 m \omega}\left(\Psi_{z},\left(\hat{a}^{2}+\hat{a}^{\dagger 2}+\hat{a}^{\dagger} \hat{a}+1+\hat{a}^{\dagger} \hat{a}\right) \Psi_{z}\right)  \\
& =\frac{\hbar}{2 m \omega}\left(z^{2}+z^{* 2}+2 z^{*} z+1\right) \\
& =\frac{\hbar}{2 m \omega}\left[\left(z+z^{*}\right)^{2}+1\right] \\
\left\langle\hat{p}^{2}\right\rangle & =-\frac{\hbar m \omega}{2}\left(\Psi_{z},\left(\hat{a}-\hat{a}^{\dagger}\right)^{2} \Psi_{z}\right) \\
& =-\frac{\hbar m \omega}{2}\left(\Psi_{z},\left(\hat{a}^{2}-\hat{a} \hat{a}^{\dagger}-\hat{a}^{\dagger} \hat{a}+\hat{a}^{\dagger 2}\right) \Psi_{z}\right) \\
& =-\frac{\hbar m \omega}{2}\left(\Psi_{z},\left(\hat{a}^{2}-\hat{a}^{\dagger} \hat{a}-1-\hat{a}^{\dagger} \hat{a}+\hat{a}^{\dagger 2}\right) \Psi_{z}\right)  \\
& =-\frac{\hbar m \omega}{2}\left(z^{2}+z^{* 2}-2 z^{*} z-1\right) \\
& =-\frac{\hbar m \omega}{2}\left[\left(z-z^{*}\right)^{2}-1\right]
\end{align}

The calculations of squared expectation values reveal the variance properties of coherent states.

\begin{align}
\Delta x & =\sqrt{\frac{\hbar}{2 m \omega}\left[\left(z+z^{*}\right)^{2}+1\right]-\frac{\hbar}{2 m \omega}\left(z+z^{*}\right)^{2}}=\sqrt{\frac{\hbar}{2 m \omega}}  \\
\Delta p & =\sqrt{-\frac{\hbar m \omega}{2}\left[\left(z-z^{*}\right)^{2}-1\right]+\frac{\hbar m \omega}{2}\left(z-z^{*}\right)^{2}}=\sqrt{\frac{\hbar m \omega}{2}}
\end{align}

Through algebraic simplification, we find that the uncertainties in position and momentum are independent of the specific coherent state parameter.

\begin{equation}
\Delta x \Delta p=\sqrt{\frac{\hbar}{2 m \omega} \frac{\hbar m \omega}{2}}=\frac{\hbar}{2}
\end{equation}

This remarkable result demonstrates that coherent states achieve the minimum possible uncertainty product allowed by quantum mechanics.

\section{Construction of coherent states via eigenfunctions}
We may alternatively build coherent states beginning with the harmonic oscillator eigenfunctions:

\begin{equation}
\frac{1}{\sqrt{2^{n} n!}} \frac{1}{\sqrt{\lambda \sqrt{\pi}}} \exp \left(-\frac{x^{2}}{2 \lambda^{2}}\right) H_{n}\left(\frac{x}{\lambda}\right)
\end{equation}

Our goal is to identify a quantum state $\Psi_{z}$ that satisfies the eigenvalue equation:

\begin{equation}
\hat{a} \Psi_{z}=z \Psi_{z}
\end{equation}

We can express this state through a Fourier expansion using the complete set of eigenfunctions $\Psi_{n}$:

\begin{equation}
\Psi_{z}=\sum_{n=0}^{\infty} c_{n} \Psi_{n} \quad \text { s.t. } \quad \hat{a} \Psi_{z}=z \Psi_{z}
\end{equation}

Applying the annihilation operator to this expansion yields:

\begin{equation}
\hat{a} \Psi_{z}=\hat{a} \sum_{n=0}^{\infty} c_{n} \Psi_{n}=\sum_{n=0}^{\infty} c_{n} \hat{a} \Psi_{n}=\sum_{n=0}^{\infty} c_{n} \sqrt{n} \Psi_{n-1}
\end{equation}

Through an index transformation $n \rightarrow n+1$, we can rewrite this as:

\begin{equation}
\hat{a} \Psi_{z}=\sum_{n=0}^{\infty} c_{n+1} \sqrt{n+1} \Psi_{n}
\end{equation}

From the eigenvalue condition, we also have:

\begin{equation}
\hat{a} \Psi_{z}=z \Psi_{z}=z \sum_{n=0}^{\infty} c_{n} \Psi_{n}
\end{equation}

\begin{equation}
z \sum_{n=0}^{\infty} c_{n} \Psi_{n}=\sum_{n=0}^{\infty} c_{n+1} \sqrt{n+1} \Psi_{n}
\end{equation}

Comparing coefficients on both sides of this equation reveals a fundamental relationship between successive terms in our expansion.

\begin{equation}
z c_{n}=\sqrt{n+1} c_{n+1} \Longrightarrow c_{n}=\frac{z c_{n-1}}{\sqrt{n}}
\end{equation}

This recursive relation provides a systematic way to determine all coefficients from the initial value $c_0$.

\begin{equation}
c_{n}=\frac{z c_{n-1}}{\sqrt{n}}=\frac{z}{\sqrt{n}} \frac{z c_{n-2}}{\sqrt{n-1}}=\ldots=\frac{z^{n}}{\sqrt{n!}} c_{0}
\end{equation}

Through iterative application of the recurrence relation, we express each coefficient in terms of the parameter $z$ and the initial coefficient $c_0$.

\begin{equation}
\left(\Psi_{z}, \Psi_{z}\right) \stackrel{!}{=} 1
\end{equation}

The normalization requirement ensures our coherent state has unit probability across all space.

\begin{align}
\left(\Psi_{z}, \Psi_{z}\right) & =\left(\sum_{n} \frac{z^{n}}{\sqrt{n!}} c_{0} \Psi_{n}, \sum_{m} \frac{z^{m}}{\sqrt{m!}} c_{0} \Psi_{m}\right)=c_{0}^{2} \sum_{n} \sum_{m} \frac{z^{n} z^{* m}}{\sqrt{n!m!}}\left(\Psi_{n}, \Psi_{m}\right)= \\
& =c_{0}^{2} \sum_{n} \sum_{m} \frac{z^{n} z^{* m}}{\sqrt{n!m!}} \delta_{n m}=c_{0}^{2} \sum_{n} \frac{\left(|z|^{2}\right)^{n}}{n!}= \\
& =c_{0}^{2} \exp \left(|z|^{2}\right)
\end{align}

The orthonormality of the basis states simplifies the double sum to a recognizable exponential series.

\begin{equation}
c_{0}^{2} \exp \left(|z|^{2}\right)=1 \Longrightarrow c_{0}=\exp \left(-\frac{|z|^{2}}{2}\right)
\end{equation}

Solving for the normalization constant yields a characteristic exponential form that ensures proper probability interpretation.

\begin{equation}
\Psi_{z}=\mathrm{e}^{-|z|^{2} / 2} \sum_{n} \frac{z^{n}}{\sqrt{n!}} \Psi_{n}
\end{equation}

The complete coherent state representation elegantly combines the basis states with coefficients determined by the parameter $z$.

\section{Temporal evolution of coherent states}
We now examine how coherent states evolve under the influence of the harmonic oscillator Hamiltonian:

\begin{align}
\exp \left(-i \frac{t}{\hbar} H\right) \Psi_{z} & =\exp \left(-|z|^{2} / 2\right) \sum_{n} \frac{z^{n}}{\sqrt{n!}} \exp \left(-i \frac{t}{\hbar} H\right) \Psi_{n}= \\
& =\exp \left(-|z|^{2} / 2\right) \sum_{n} \frac{z^{n}}{\sqrt{n!}} \exp \left(-i \frac{t}{\hbar} \hbar \omega\left(\hat{n}+\frac{1}{2}\right)\right) \Psi_{n}= \\
& =\exp \left(-|z|^{2} / 2\right) \sum_{n} \frac{z^{n}}{\sqrt{n!}} \exp \left(-i t \omega\left(n+\frac{1}{2}\right)\right) \Psi_{n}= \\
& =\exp \left(-\frac{|z|^{2}+i \omega t}{2}\right) \sum_{n} \frac{z^{n}}{\sqrt{n!}}[\exp (-i t \omega)]^{n} \Psi_{n}= \\
& =\exp \left(-\frac{|z|^{2}+i \omega t}{2}\right) \sum_{n} \frac{z^{n}(t)}{\sqrt{n!}} \Psi_{n}=
\end{align}

The time evolution maintains the coherent state structure with a time-dependent parameter $z(t)=z \exp (-i t \omega)$.

We observe that the expectation values follow classical-like trajectories:

$\langle x\rangle=\sqrt{\frac{\hbar}{2 m \omega}}\left(\Psi_{z},\left(a^{\dagger}+a\right) \Psi_{z}\right)=\sqrt{\frac{\hbar}{2 m \omega}}\left(z^{*}(t)+z(t)\right)=\sqrt{\frac{2 \hbar}{m \omega}}\left|z_{0}\right| \cos \left(\omega t-\phi_{0}\right)$

\begin{equation}
\langle p\rangle=\sqrt{\frac{\hbar m \omega}{2}}\left(\Phi_{z},\left(a-a^{\dagger}\right) \Phi_{z}\right)=\sqrt{\frac{\hbar m \omega}{2}}\left(z^{*}(t)-z(t)\right)=\sqrt{\frac{\hbar m \omega}{2}}\left|z_{0}\right| \sin \left(\omega t-\phi_{0}\right)
\end{equation}

These expressions demonstrate that coherent states exhibit dynamics remarkably similar to classical harmonic oscillation with amplitude $A=\sqrt{\frac{2 \hbar}{m \omega}}\left|z_{0}\right|$.

\section{Complementary approach to coherent state dynamics}
An alternative derivation of coherent state evolution begins with the ansatz:

\begin{equation}
\Psi_{z}(x, t)=\mathrm{e}^{i \alpha} \Psi_{z}(x)
\end{equation}

Where the parameter $z$ evolves according to:

\begin{equation}
z(t)=z_{0} \mathrm{e}^{-i \omega t}
\end{equation}

\begin{align}
i \hbar \frac{\partial}{\partial t}\left(\mathrm{e}^{i \alpha} \Psi_{z}(x)\right) & =\mathrm{e}^{i \alpha} H \Psi_{z}(x) \\
i \hbar \mathrm{e}^{i \alpha} \frac{\partial \Psi_{z}(x)}{\partial t}-\hbar \dot{\alpha} \mathrm{e}^{i \alpha} \Psi_{z}(x) & =\mathrm{e}^{i \alpha} H \Psi_{z}(x)
\end{align}

The Schrödinger equation governs how our quantum state evolves in time, capturing both the phase evolution and the change in state structure.

\begin{align}
& -\hbar \dot{\alpha} \mathrm{e}^{i \alpha} \Psi_{z}(x)+i \hbar \mathrm{e}^{i \alpha} \frac{\partial}{\partial t} \sum_{n=0}^{\infty} \frac{z^{n}}{\sqrt{n!}} \Psi_{n}= \\
& =-\hbar \dot{\alpha} \mathrm{e}^{i \alpha} \Psi_{z}(x)+i \hbar \mathrm{e}^{i \alpha} \sum_{n=0}^{\infty} n \frac{\dot{z}}{z} \frac{z^{n}}{\sqrt{n!}} \Psi_{n}= \\
& =-\hbar \dot{\alpha} \mathrm{e}^{i \alpha} \Psi_{z}(x)+i \hbar \mathrm{e}^{i \alpha} \sum_{n=0}^{\infty} n \frac{-i \omega \cancel{z}}{\cancel{z}} \frac{z^{n}}{\sqrt{n!}} \Psi_{n}=  \\
& =-\hbar \dot{\alpha} \mathrm{e}^{i \alpha} \Psi_{z}(x)+\mathrm{e}^{i \alpha} \sum_{n=0}^{\infty}(\hbar \omega n) \frac{z^{n}}{\sqrt{n!}} \Psi_{n}=
\end{align}

Substituting our time-dependent parameter into the coherent state expansion reveals how each energy component evolves.

\begin{align}
& \mathrm{e}^{i \alpha} H \Psi_{z}(x)=\mathrm{e}^{i \alpha} \hbar \omega\left(\hat{n}+\frac{1}{2}\right) \sum_{n=0}^{\infty} \frac{z^{n}}{\sqrt{n!}} \Psi_{n}=  \\
& =\mathrm{e}^{i \alpha} \sum_{n=0}^{\infty} \hbar \omega n \frac{z^{n}}{\sqrt{n!}} \Psi_{n}+\frac{\hbar \omega}{2} \mathrm{e}^{i \alpha} \Psi_{z}
\end{align}

The Hamiltonian action on our coherent state separates into terms representing excited states and the zero-point energy contribution.

\begin{align}
-\hbar \dot{\alpha} e^{i \alpha} \Psi_{z}(x)+e^{i \alpha} \sum_{n=0}^{\infty}(\hbar \omega n) \frac{z^{n}}{\sqrt{n!}} \Psi_{n} & =e^{i \alpha} \sum_{n=0}^{\infty} \hbar \omega n \frac{z^{n}}{\sqrt{n!}} \Psi_{n}+\frac{\hbar \omega}{2} e^{i \alpha} \Psi_{z}  \\
-\hbar \dot{\alpha} & =\frac{\hbar \omega}{2}
\end{align}

Comparing terms in the Schrödinger equation determines the phase evolution of our coherent state.

\begin{equation}
\alpha=-\frac{\omega}{2} t
\end{equation}

The phase accumulates linearly with time, reflecting the zero-point energy contribution to the coherent state dynamics.
